===== BEGIN FILE: macros.tex =====
\newtheorem{theorem}{Theorem}[section]
\newtheorem{lemma}[theorem]{Lemma}
\newtheorem{corollary}[theorem]{Corollary}
\theoremstyle{definition}
\newtheorem{definition}[theorem]{Definition}
\newtheorem{proposition}[theorem]{Proposition}
\theoremstyle{remark}
\newtheorem*{remark}{Remark}








\newcommand{\calA}{\mathcal{A}}
\newcommand{\calB}{\mathcal{B}}
\newcommand{\calC}{\mathcal{C}}
\newcommand{\calD}{\mathcal{D}}
\newcommand{\calE}{\mathcal{E}}
\newcommand{\calF}{\mathcal{F}}
\newcommand{\calG}{\mathcal{G}}
\newcommand{\calH}{\mathcal{H}}
\newcommand{\calI}{\mathcal{I}}
\newcommand{\calJ}{\mathcal{J}}
\newcommand{\calK}{\mathcal{K}}
\newcommand{\calL}{\mathcal{L}}
\newcommand{\calM}{\mathcal{M}}
\newcommand{\calN}{\mathcal{N}}
\newcommand{\calO}{\mathcal{O}}
\newcommand{\calP}{\mathcal{P}}
\newcommand{\calQ}{\mathcal{Q}}
\newcommand{\calR}{\mathcal{R}}
\newcommand{\calS}{\mathcal{S}}
\newcommand{\calT}{\mathcal{T}}
\newcommand{\calU}{\mathcal{U}}
\newcommand{\calV}{\mathcal{V}}
\newcommand{\calW}{\mathcal{W}}
\newcommand{\calX}{\mathcal{X}}
\newcommand{\calY}{\mathcal{Y}}
\newcommand{\calZ}{\mathcal{Z}}

\newcommand{\e}{\mathrm{e}}
\renewcommand{\i}{\mathrm{i}}
\renewcommand{\d}{\mathrm{d}}
\newcommand{\id}{\mathrm{id}}
\newcommand{\im}{\mathop{\mathrm{im}}}
\newcommand{\re}{\mathop{\mathrm{re}}}
\newcommand{\sgn}{\mathop{\mathrm{sgn}}}
\renewcommand{\vec}[1]{\boldsymbol{#1}}
\newcommand{\tr}{\mathrm{tr}}
\newcommand{\Z}{\mathbb{Z}}
\newcommand{\Q}{\mathbb{Q}}
\newcommand{\R}{\mathbb{R}}
\newcommand{\C}{\mathbb{C}}



\newcommand{\Hom}{\mathrm{Hom}}
\newcommand{\Ext}{\mathrm{Ext}}
\newcommand{\Tor}{\mathrm{Tor}}
\newcommand{\Aut}{\mathrm{Aut}}
\newcommand{\End}{\mathrm{End}}
\newcommand{\tH}{\widetilde{H}}
\newcommand{\tOmega}{\widetilde{\Omega}}
\newcommand{\ch}{\mathrm{ch}}
\newcommand{\Sq}{\mathrm{Sq}}
\newcommand{\Arf}{\mathrm{Arf}}

\newcommand{\Rep}{\mathsf{Rep}}
\newcommand{\sRep}{\mathsf{sRep}}
\newcommand{\Vect}{\mathsf{Vec}}
\newcommand{\TY}{\mathsf{TY}}
\newcommand{\Fun}{\mathsf{Fun}}
\newcommand{\Obj}{\mathsf{Obj}}
\newcommand{\Mor}{\mathsf{Hom}}

\newcommand{\Det}{\calD et\,}
\newcommand{\Pf}{\calP\!f\,}
\newcommand{\G}{\mathcal{G}}
\newcommand{\U}{\mathcal{U}}
\renewcommand{\H}{\mathbb{H}}
\newcommand{\T}{\mathcal{T}}
\usepackage{stackengine}
\newcommand\oast{\stackMath\mathbin{\stackinset{c}{0ex}{c}{0ex}{\ast}{\bigcirc}}}
\def\vev#1{{\langle{#1}\rangle}}







\def\sc#1{\textcolor{blue}{{\bf SC}: #1}}
\def\ac#1{\textcolor{red}{{\bf AC}: #1}}
\def\rp#1{\textcolor{purple}{{\bf RP}: #1}}
\def\mn#1{\textcolor{magenta}{{\bf MN}: #1}}
\def\sb#1{\textcolor{cyan}{[{\bf SB}: #1]}}


\usepackage{titlesec}


\newcounter{subsubsubsection}[subsubsection]
\renewcommand\thesubsubsubsection{\thesubsubsection.\arabic{subsubsubsection}}


\newcommand{\subsubsubsection}[1]{
  \refstepcounter{subsubsubsection}
  \addcontentsline{toc}{paragraph}{\protect\numberline{\thesubsubsubsection}#1}
  \paragraph{\thesubsubsubsection\quad #1}
}


\titleformat{\paragraph}[block]{\normalfont\normalsize\bfseries}{\theparagraph}{1em}{}


\titlespacing*{\paragraph}{0pt}{1.5ex plus 0.5ex minus .2ex}{0.8ex plus .2ex}
===== END FILE: macros.tex =====

===== BEGIN FILE: main.tex =====
\documentclass[
 reprint,







nofootinbib,


 amsmath,amssymb,
 aps,
 prd






tightenlines
]{revtex4-2}

\usepackage{graphicx}
\usepackage{dcolumn}
\usepackage{bm}
\usepackage[hidelinks,colorlinks=true,allcolors=blue]{hyperref}
\usepackage{orcidlink}

\usepackage{color}
\usepackage{slashed}
\usepackage[utf8]{inputenc}
\usepackage{amsmath,amssymb,multirow,graphicx, subfigure}


\usepackage{amsthm}
\usepackage{braket}
\usepackage[T1]{fontenc}
\usepackage{enumitem}

\input{macros.tex}

\begin{document}

\title{
Does hot QCD have a conformal manifold in the chiral limit?
}

\author{Shi Chen \orcidlink{0000-0002-8554-5098}}
\email{s.chern.phys@gmail.com}
\affiliation{School of Physics and Astronomy, University of Minnesota, Minneapolis, MN 55455, USA}
\author{Aleksey Cherman \orcidlink{0000-0002-1039-8476}}
\email{acherman@umn.edu}
\affiliation{School of Physics and Astronomy, University of Minnesota, Minneapolis, MN 55455, USA}
\author{Robert D. Pisarski \orcidlink{0000-0002-7862-4759}}
\email{pisarski@bnl.gov}
\affiliation{Physics Department, Brookhaven National Laboratory,
Upton, NY 11973-5000}

\date{\today}

\begin{abstract}
Recent lattice evidence suggests the chiral phase transition in QCD is second-order for $N_f \ge 2$ massless flavors. We
constrain CFT descriptions of a critical line in
temperature $T$ and imaginary baryon chemical potential $\theta_B =
i\mu_B/T$. An 't~Hooft anomaly at general $\theta_B$ constrains the
transition even at $\theta_B = 0$, leaving only three minimal scenarios. The
best-motivated scenario for $N_f\ge3$, and perhaps also $N_f = 2$, is beyond
Ginzburg-Landau, featuring a conformal manifold of
$\theta_B$-dependent universality classes with an exactly marginal
operator related to baryon density.
\end{abstract}

\maketitle



{\bf Introduction.}


Understanding the phase diagram of quantum chromodynamics (QCD) as a
function of temperature $T$ and baryon chemical potential $\mu_B$ is a
central challenge in high-energy physics,
with implications for heavy-ion collisions, astrophysics, and the early
universe.
At $\mu_B=0$, increasing $T$ restores the spontaneously broken approximate
chiral symmetry through a crossover at $T \sim 155\,
\textrm{MeV}$~\cite{HotQCD:2018pds,Borsanyi:2020fev}.
If $N_f$ quarks are massless, chiral symmetry becomes exact, and the crossover
becomes a genuine phase transition. The standard approach to this
nonperturbative phase transition appeals to Landau universality. A
Ginzburg-Landau analysis based on the chiral condensate $\Phi = \bar{\psi}_L
\psi_R$~\cite{Pisarski:1983ms} suggests a first-order transition for $N_f\geq3$,
and a second-order transition for $N_f\!=\!2$,
in the $O(4)$ Wilson-Fisher universality class.

Early lattice simulations seemed to agree with these
predictions~\cite{Brown:1990ev,Karsch:1994hm,Iwasaki:1995ij,JLQCD:1998qth,
Karsch:2001nf,Liao:2001en,deForcrand:2003vyj,Jin:2014hea}, but recent lattice
simulations~\cite{Cuteri:2021ikv,Dini:2021hug,Nakamura:2022abk,Zhang:2022kzb,Zhang:2024ldl,Zhang:2025vns,Klinger:2025xxb,DAmbrosio:2025ldv,Klinger:2026pbe} attribute
the earlier $N_f \ge 3$ results to discretization artifacts and find no
evidence of a first-order transition for any $N_f \ge 2$ at $\mu_B=0$ or
imaginary $\mu_B$~\cite{DAmbrosio:2025ldv}. Lattice results thus suggest
that the chiral transition is second-order, or at most weakly
first-order, from $N_f = 2$ up to the onset of the conformal window\cite{Klinger:2026pbe}, at $\mu_B = 0$ and along a line
extending to imaginary $\mu_B$.

Various theoretical interpretations of the recent lattice results have been proposed~\cite{Fejos:2022mso,Bernhardt:2023hpr,Fejos:2024bgl,Giacosa:2024orp,Pisarski:2024esv}, but compelling solutions remain elusive.
Here we systematically constrain the 3D conformal field theories (CFTs) that can appear on a
natural critical line of chiral transitions in the $\theta_B$-$T$ phase
diagram, with real $\theta_B= i\mu_B / T$. We find only three minimal
scenarios compatible with the ’t Hooft
anomaly~\cite{Yonekura:2019vyz,Nishimura:2019umw,Kobayashi:2023ajk},
high-$T$ perturbative calculations~\cite{Roberge:1986mm}, and low-$T$
QCD phenomenology, shown in Fig.~\ref{fig:allowed_scenarios}.

\begin{figure*}[ht]
  \centering
  \includegraphics[width=0.3\textwidth]{figures/all_broken_standalone.pdf}
  \hfill
  \includegraphics[width=0.3\textwidth]{figures/dqcp_standalone.pdf}
  \hfill
  \includegraphics[width=0.3\textwidth]{figures/conformal_manifold_standalone.pdf}
  \caption{Minimal phase diagrams of massless QCD consistent with a natural second-order
   line of chiral phase transitions. \textbf{Left}: Landau scenario.
   \textbf{Middle}: Landau-DQCP scenario. \textbf{Right}: Conformal-manifold
   scenario. The conformal manifold scenario appears plausible for all
   $N_f \ge 2$, while the Landau and Landau-DQCP scenarios are only
   plausible for $N_f=2$.}
  \label{fig:allowed_scenarios}
\end{figure*}

In both the Landau (left) and Landau-DQCP (middle) scenarios, the critical
line flows to a single Ginzburg-Landau CFT for generic $\theta_B$, but the
latter also features a deconfined quantum critical point
(DQCP)~\cite{Senthil:2003eed,Senthil:2004fuw,Senthil:2023vqd} at $\theta_B =
\pi$. These scenarios are conceivable for $N_f=2$, but for $N_f \ge 3$ they
require a natural fixed point in the Ginzburg-Landau description of
Ref.~\cite{Pisarski:1983ms} that has not been found despite extensive
searches. In the conformal-manifold scenario (right), the critical line
is a continuum of distinct beyond-Landau 3D CFTs. We argue
that this scenario is supported by a $(2+\epsilon)$D expansion and existing
conformal bootstrap bounds for $N_f \ge 2$. So the answer to the title's
question could be \emph{yes}.



{\bf Symmetry and anomaly.}


We work in Euclidean signature throughout.
Massless QCD, a 4D $SU(N_c)$ gauge theory for $N_c\ge3$ with $N_f$ massless fundamental Dirac fermions, has the continuous global symmetry
\begin{align}
\begin{split}
    G = &\:\frac{U(1)_B \times SU(N_f)_L \times SU(N_f)_R }{\Z_{N_f}},
\end{split}
\label{eq:QCD_symmetry}
\end{align}
where $\Z_{N_f}$ sets $(\e^{2\pi i N_c/N_f}, \e^{-2\pi i/N_f}
\mathbf{1}_{N_f}, \e^{-2\pi i/N_f} \mathbf{1}_{N_f})$ to  the identity.
Here $U(1)_B=U(1)_V/\Z_{N_c}$ is the baryon symmetry, while
$SU(N_f)_{L} \times SU(N_f)_R$ is the chiral symmetry.
The $U(1)_A$ symmetry of the classical theory is destroyed by the axial anomaly, and so it is not part of $G$.
The discrete global symmetry acts non-trivially on $G$ and is generated by charge conjugation $\calC$ and geometric reflections, including time reversal.

The symmetry $G$ has an 't Hooft anomaly~\cite{tHooft:1979rat}. Turning
on background gauge fields $A_B$, $A_L$, $A_R$ for $U(1)_B$,
$SU(N_f)_L$, $SU(N_f)_R$, respectively, leads to an anomalous
background gauge variation described by a 5D invertible field
theory~\cite{Freed:2014iua,Kapustin:2014lwa,Witten:2015aba,Yonekura:2016wuc,Bhardwaj:2023kri}
with partition function
\begin{align}\label{eq:SPT_4d}
\begin{gathered}
    \mathcal{I}_{5D}(A_B,A_L,A_R) = \\
    \mathcal{J}_{5D}^{N_c}(A_L)\mathcal{J}_{5D}^{-N_c}(A_R)
    \exp\!\left[\frac{i}{8\pi^2} \!\int_{5D}\!\! A_B \tr( F_L^2 - F_R^2)\right]\!,
\end{gathered}
\end{align}
where $F_{L,R}$ is $A_{L,R}$'s field strength.
For $N_f \ge 3$, $\calJ_{5D}(A)$ is the 5D Chern-Simons partition
function~\cite{Wess:1971yu}, while for $N_f=2$ it is the 5D spin
topological invariant associated with the Witten
anomaly~\cite{Witten:1982fp,Wang:2018qoy}.

As a renormalization group (RG) invariant, the 't Hooft anomaly~\eqref{eq:SPT_4d} ensures massless excitations~\cite{Frishman:1980dq,Coleman:1982yg}, but their nature and interactions depend on $N_c$ and $N_f$.
Below the conformal window (very roughly
$N_f\lesssim\frac{8}{3}N_c$~\cite{LatticeStrongDynamics:2018hun,Kuti:2022ldb,Ingoldby:2023mtf,Hasenfratz:2023wbr,Hasenfratz:2024fad,Bergner:2025yke}),
the continuous symmetry $G$ is spontaneously broken to
\begin{equation}
    H = \frac{U(1)_B\times SU(N_f)_V}{\Z_{N_f}}\,,
\end{equation}
while the discrete symmetries remain intact.
The gapless IR physics is described via a nonlinear scalar field
\begin{equation}\label{eq:U(x)}
    \phi(x) \ \in\ G/H\ \simeq\ SU(N_f)\,.
\end{equation}
In this phase, $\mathcal{J}^{N_c}(A_L)\mathcal{J}^{-N_c}(A_R)$ in Eq.~\eqref{eq:SPT_4d} is matched by the Wess-Zumino-Witten action~\cite{Witten:1983tw} for $N_f\ge3$, and by the discrete spin $\theta$-angle~\cite{Witten:1983tw} for $N_f=2$.
The $A_B$ term is matched by identifying the Skyrme current (from $\pi_3(SU(N_f)) = \Z$) with the baryon current~\cite{Witten:1983tx}.



{\bf Circle compactification.}


We now put massless QCD on $\R^3\times S^1$ with circumference $\beta$.
An $S^1$-homogeneous gauge field splits into a gauge field on $\R^3$
and a holonomy scalar field on $\R^3$. In particular, the $U(1)_B$
holonomy gives a $2\pi$-periodic scalar field $\theta_B$.

\footnote{What directly shows up in the QCD Lagrangian is $U(1)_V$ instead of $U(1)_B$. Thus in
the path integral, we have to define $\theta_B$ via $\theta_B =\theta_VN_c$.
Then the shift
$\theta_B\sim\theta_B+2\pi$ is generated by an improper $SU(N_c)$ gauge transformation $g$ that
satisfies $g(x,\tau+\beta) = g(x,\tau)\,\e^{2\pi i/N_c}$ for $\tau\in
S^1$~\cite{Roberge:1986mm}.}

For constant $\theta_B$, the circle compactification has two Hamiltonian
interpretations, and all physical observables, including the free energy
density
\begin{equation}
    \label{eq:free_energy}
    f_{\beta}(\theta_B)\equiv-\lim_{V\to\infty}\frac{1}{\beta V}\ln Z_{\beta}(\theta_B),
\end{equation}
must be \emph{simultaneously} interpretable in both pictures:
\begin{itemize}
    \item \textbf{(Thermal)} Hamiltonian evolution along $S^1$ yields a grand
    canonical ensemble on $\R^3$ with temperature $T=1/\beta$ and baryonic chemical potential $\mu_B =
    -iT\theta_B$. Then $-f_{\beta}(\theta_B)$ is the thermal pressure.
    \item \textbf{(Quantum)} Hamiltonian evolution along $\R$ yields a
    zero-temperature quantum system on $\R^2\!\times\!S^1$ with
    circumference $\beta$ and a baryonic Aharonov-Bohm phase
    $e^{i\theta_B}$. Then $f_{\beta}(\theta_B)$ is the Casimir energy
    density.
\end{itemize}
Therefore, with real $\theta_B\neq0$, we lose reflection positivity
reflecting $S^1$, but preserve reflection positivity reflecting an
$\R \subset \R^3$.
The imaginary baryon density at zero Matsubara frequency,
\begin{equation}\label{eq:baryon_density}
    \rho_B(x) \equiv \frac{i}{\beta}\int_{0}^{\beta}\!\!\!\d\tau\, j_B^0(x,\tau)\,,
\end{equation}
is anti-hermitian in the thermal picture, but is hermitian in the quantum picture.

The quantum picture is convenient for describing 3D infrared (IR)
effective theories at length scales $\ell \gg\beta$. These theories are
unitary 3D quantum field theories (QFTs) for any real $\theta_B$, and $\rho_B$ flows to a hermitian operator thereof. When the 3D effective
theory becomes scale-invariant at a critical point,
unitarity is expected to enhance scale invariance to full conformal invariance; see
e.g.~Refs.~\cite{Polchinski:1987dy,El-Showk:2011xbs,Nakayama:2013is,Delamotte:2015aaa,Paulos:2015jfa,Meneses:2018xpu}.

Upon $S^1$ compactification, the 5D invertible theory~\eqref{eq:SPT_4d}
reduces to the 4D invertible theory
\begin{align}\label{eq:SPT_3D}
    \mathcal{I}_{4D}(\theta_B,A_L,A_R) =
    \exp\!\left[\frac{i}{8\pi^2}\! \int_{4D}\!\! \theta_B \,\tr (F_L^2 -F_R^2)\right]\!.
\end{align}
It is well-defined on closed 4D oriented manifolds, but is not
invariant under $\theta_B\to\theta_B+2\pi$ if the 4D manifold has a 3D
boundary. Hence, with $S^1$-homogeneous $A_{L,R}$, the partition
function of $S^1$-compactified massless QCD
satisfies~\cite{Yonekura:2019vyz,Nishimura:2019umw,Kobayashi:2023ajk}
\begin{align}\label{eq:anomaly_3D}
\begin{gathered}
    Z_{\beta}(\theta_B + 2\pi, A_L, A_R) \\
    = Z_{\beta}(\theta_B, A_L, A_R) \,\calJ_{3D}(A_L) \,\calJ_{3D}^{-1}(A_R)
\end{gathered}
\end{align}
where $\calJ_{3D}(A)$ is the 3D Chern-Simons partition function
\begin{align}
    \calJ_{3D}(A) =  \exp\left[\frac{\i}{4\pi}\int_{3D}
    \mathrm{tr}\left( A \d A - \i\frac{2}{3} A^3 \right)\right] \,.
\end{align}
This is a mixed anomaly between $SU(N_f)_{L,R}$ and a $2\pi$-periodic coupling~\cite{Kikuchi:2017pcp,Cordova:2019jnf,Cordova:2019uob,Heidenreich:2021xpr,Garcia-Valdecasas:2024cqn,Yu:2024jtk}, or a ``$(-1)$-form $U(1)$ symmetry''~\cite{Najjar:2024vmm,Aloni:2024jpb, Lin:2025oml,Robbins:2025urk}.

Upon $S^1$ compactification, the discrete symmetry is generated by $\calC$, the $S^1$-reflection $\calS$, and $\R$-reflections $\calR$ for $\R\subset\R^3$.
They act on $S^1$-homogeneous 3D background fields as
\begin{equation}
\!\!\!\!\begin{gathered}
    \theta_B(x) \xrightarrow{\calC} -\theta_B(x),\quad A_{L,R}(x) \xrightarrow{\calC} -A_{R,L}^*(x),\\
    \theta_B(x) \xrightarrow{\calS} -\theta_B(x),\quad A_{L,R}(x) \xrightarrow{\calS} A_{R,L}(x),\\
    \theta_B(x) \xrightarrow{\calR}\theta_B(\overline{x}),\quad A_{L,R}(x) \xrightarrow{\calR} \overline{A}_{R,L}(\overline{x}),
\end{gathered}
\end{equation}
where the bars indicate an $\R$-reflection.
$\calC$, $\calS$, and $\calR$ exchange $L$ and $R$.
When $N_f=2$, $\calC\calS$ can be undone by $SU(N_f)_{L,R}$.

A constant $\theta_B\neq0,\pi$ preserves only the subgroup generated by
$\calC\calS$ and $\calR$, whereas $\calS$, $\calC$, $\calS\calR$, and
$\calC\calR$ are restored at $\theta_B=\pi$ by the $2\pi$-periodicity
of $\theta_B$. Therefore, at $\theta_B=\pi$,
Eq.~\eqref{eq:anomaly_3D} gives a mixed anomaly between $SU(N_f)_{L,R}$
and $(\Z_2)_{\calS\calR}$~\cite{Yonekura:2019vyz,Nishimura:2019umw,Kobayashi:2023ajk}.
Equivalently, this is an anomaly of
$SU(N_f)_{L,R}\rtimes(\Z_2)_{\calC\calR}$, or of
$[SU(N_f)_L\times SU(N_f)_R]\rtimes(\Z_2)_{\calS}$ or
$\rtimes(\Z_2)_{\calC}$. The 4D invertible theory~\eqref{eq:SPT_3D} captures
this because, at $\theta_B=\pi$, Eq.~\eqref{eq:SPT_3D} can be defined on
unorientable manifolds or with twists that exchange $L$ and $R$.



{\bf Low-$T$ and high-$T$ Phases.}


Although it does not guarantee gapless modes, the 3D
anomaly~\eqref{eq:SPT_3D} excludes $\theta_B$-uniform trivially gapped
phases, as well as a trivially gapped phase at $\theta_B=\pi$.
The 3D IR physics at length scales $\ell\gg\beta$ must see $\theta_B$ and the symmetry
\begin{equation}\label{eq:3D_IR_symmetry}
    \frac{SU(N_f)_L\times SU(N_f)_R}{\Z_{N_f}}
    \begin{cases}
        \!\rtimes (\Z_2)^3_{\calC,\calS,\calR}, & \theta_B=0,\pi \\
        \!\rtimes (\Z_2)^2_{\calC\calS,\calR}, & \theta_B\neq0,\pi
    \end{cases}
\end{equation}
where discrete symmetry needs to quotient out $(\Z_2)_{\calC\calS}$ for $N_f=2$.
The operator $\rho_B$ of Eq.~\eqref{eq:baryon_density} is neutral under
the second line of Eq.~\eqref{eq:3D_IR_symmetry} but charged under the
first.
In contrast, the 3D baryon symmetry $U(1)_B$ and chiral holonomies $\theta_{L,R}$ need not appear in the IR.

At sufficiently low $T$, the symmetry $G$ is expected to remain broken to
$H$ just as at $T=0$, and the IR physics is described by the 3D
reduction of the nonlinear scalar field $\phi(x)$ of
Eq.~\eqref{eq:U(x)}.
The 3D symmetry~\eqref{eq:3D_IR_symmetry} acts on it by $\phi(x)\to L\phi(x)R^{-1}$ for $L,R\in SU(N_f)$ and
\begin{equation}
    \phi(x) \xrightarrow{\calC} \phi^T(x)\,,\
    \phi(x) \xrightarrow{\calS} \phi^{\dag}(x)\,,\
    \phi(x) \xrightarrow{\calR} \phi^{\dag}(\overline{x})\,.
\end{equation}
When $N_f=2$, $\calC\calS$ is absorbed into $SU(2)_{L,R}$.
The low-$T$ 3D IR effective QFT is given by~\cite{Xu:2011sj,Yonekura:2019vyz}
\begin{equation}\label{eq:lowT_EFT}
\begin{gathered}
    \!\!\!\int\!\!\calD \phi\exp\!\left\{\!-\frac{1}{g^2}\!\!\int\!\!\d^3\!x\,\tr \partial_{a}\phi\, \partial^{a}\!\phi^{\dagger}\! + i\!\!\int\!\!\theta_B\frac{\tr(\phi^{\dagger}\d \phi)^3}{24\pi^2} \!\right\}\!.
\end{gathered}
\end{equation}
To verify anomaly matching in Eq.~\eqref{eq:lowT_EFT}, one can turn on background
gauge fields $A_{L,R}$ and study their interplay with the Skyrme 3-form~\cite{Witten:1983tw,Xu:2011sj,Yonekura:2019vyz}.
In this low-$T$ gapless phase, varying $\theta_B$ does not induce
phase transitions since baryons are heavy.

At sufficiently high $T$, perturbative evaluations of some quantities
are trustworthy; see e.g.~Refs.~\cite{Gross:1980br,KapustaGale201102}.
They imply that chiral symmetry is restored and there is a mass gap for
all $\theta_B$~\cite{Roberge:1986mm}. Roberge
and Weiss~\cite{Roberge:1986mm} pointed out a first-order phase transition at $\theta_B=\pi$.
This is associated with discrete symmetry breaking at $\theta_B=\pi$,
\begin{equation}
\label{eq:broken_Z2}
    (\Z_2)^3_{\calC,\calS,\calR} \ \ \to\ \ (\Z_2)^2_{\calC\calS,\calR},
\end{equation}
which needs to quotient out $(\Z_2)_{\calC\calS}$ for $N_f=2$.
The order operator is $\rho_B$ with
\begin{equation}
    \langle\rho_B\rangle = \frac{\partial f_{\beta}(\theta_B)}{\partial\theta_B}\,.
\end{equation}
In this high-$T$ gapped phase, continuous chiral symmetry flows to an emergent $2$-form $U(1)$ symmetry~\cite{Gaiotto:2014kfa} and the 3D anomaly~\eqref{eq:SPT_3D} is matched via symmetry transmutation~\cite{Seiberg:2025bqy}; we explain this high-$T$ IR effective theory in the Supplemental Material~\ref{sec:high_T_EFT}.

In summary, chiral symmetry is broken in the low-$T$ phase and restored
in the high-$T$ phase, while the discrete symmetry at $\theta_B=\pi$ is
broken

\footnote{
If a CFT matching the symmetry and anomaly existed, a second-order transition at $\theta_B=\pi$ without symmetry breaking could occur at lower $T$ in the high-$T$ phase~\cite{Xu:2011sj,Yonekura:2019vyz,Nishimura:2019umw,Kobayashi:2023ajk}; see the Supplemental Material~\ref{sec:anomalons} for comments.}

in the high-$T$ phase and restored in the low-$T$ phase. This
complementarity of symmetry realizations is a fingerprint of mixed
anomalies.



{\bf Consistent phase diagrams.}


The low-$T$ and high-$T$ phases are separated by a critical line
$\theta_B \mapsto T_c(\theta_B)$. Motivated by recent lattice
results~\cite{Cuteri:2021ikv,Dini:2021hug,Nakamura:2022abk,Zhang:2022kzb,Zhang:2024ldl,Zhang:2025vns,Klinger:2025xxb,DAmbrosio:2025ldv,Klinger:2026pbe},
we postulate that this line generically consists of natural second-order
points with emergent scale invariance.
The IR physics of such a natural critical point is a 3D unitary CFT whose neutral
sector under the symmetry~\eqref{eq:3D_IR_symmetry} has precisely one
relevant primary scalar operator $\calO_{E}$, to which the $T$
perturbation flows. The operator $\rho_B$ also flows to
$\calO_{E}$ at $\theta_B \neq 0,\pi$, or at $\theta_B = 0,\pi$ if
spontaneous symmetry breaking~\eqref{eq:broken_Z2} occurs.

The 't Hooft anomaly~\eqref{eq:SPT_3D} prohibits the critical line from
flowing to a single CFT for all $\theta_B$. First, to produce
the behavior~\eqref{eq:anomaly_3D}, the IR physics must change as we
smoothly vary $\theta_B$ to $\theta_B+2\pi$. Second, it must differ at
$\theta_B=0$ and $\pi$, since these points have the same symmetry but
distinct anomalies. Therefore, the IR physics must change from
$\theta_B=0$ to $\theta_B=\pi$, either abruptly or gradually. Abrupt
changes correspond to phase transitions at special $\theta_B$, like in
the high-$T$ phase. Gradual changes, as in the low-$T$ phase,
correspond to smooth variations of the CFT data along the critical line.

If the changes in the IR physics are abrupt, a neighborhood of
$\theta_B=0$ on the critical line only cares about the breaking
$G \to H$, and should fit the Landau paradigm, i.e.~be described by a
natural fixed point of the Ginzburg-Landau model of
Ref.~\cite{Pisarski:1983ms}. Minimally, it suffices to have one
phase transition at $\theta_B = \pi$ and a
Ginzburg-Landau description for the rest of the critical line.
If the transition at $\theta_B = \pi$ is first-order ({\bf Landau
scenario}, Fig.~\ref{fig:allowed_scenarios} left), discrete
symmetry breaking~\eqref{eq:broken_Z2} with $\langle\rho_B\rangle\neq0$ occurs at $\theta_B = \pi$. There must then be a
window at $\theta_B=\pi$ where chiral and discrete symmetries are
broken, with an Ising transition at its lower edge.

\footnote{Such a window is disfavored by current lattice evidence~\cite{Bonati:2018fvg,Cuteri:2022vwk}.}
If instead the
transition at $\theta_B = \pi$ is second order ({\bf Landau-DQCP
scenario}, Fig.~\ref{fig:allowed_scenarios} middle), it must be a
beyond-Landau critical point analogous to condensed-matter
DQCPs~\cite{Senthil:2003eed,Senthil:2004fuw,Senthil:2023vqd}. The
operator $\rho_B$ flows to a relevant (or marginally relevant) operator
at this DQCP.

If the changes in the IR physics are gradual, the CFT data instead vary smoothly with $\theta_B$ on the critical line ({\bf conformal-manifold scenario}, Fig.~\ref{fig:allowed_scenarios} right).
This scenario requires a conformal manifold --- a continuum of distinct
CFTs --- generated by an exactly marginal operator $\calO_{B}$. The
operator $\rho_B$ flows to $\calO_{B}$ at $\theta_B=0,\pi$, and to a
linear combination of $\calO_{B}$ and $\calO_{E}$ at
$\theta_B\neq0,\pi$. The presence of $\calO_{B}$ makes this scenario
intrinsically beyond Ginzburg-Landau.

One can imagine more complicated scenarios with phase transitions
away from $\theta_B=\pi$, but they generically require unnatural
multicritical points. In any case, the three scenarios above are the
minimal ones under the assumption of a second-order critical line.



{\bf Landau scenarios.}


The Landau and Landau-DQCP scenarios posit that the critical line at
$\theta_B \neq \pi$ is described

\footnote{
For $\theta_B\neq0$, we should allow chiral-symmetric $(\Z_2)^2_{\calC\calS,\calR}$-neutral $(\Z_2)^3_{\calC,\calS,\calR}$-charged deformations, such as $i\lambda_4 \int[\tr (\Phi^\dagger d \Phi)^3 + \text{c.c.}]$.
However, operators in this sector are expected to be irrelevant.}

by a natural fixed point of the 3D linear sigma model of
Ref.~\cite{Pisarski:1983ms}:
\begin{align}\label{eq:Landau-Ginzburg}
\begin{gathered}
    \int\!\calD\Phi\,\exp\Big\{\!-\!\int\!\d^3x\Big[\tr \partial_{a}\Phi\,\partial^{a}\Phi^{\dag} + \mu\,\tr\Phi\Phi^{\dag} \\
    +\, \lambda_1\!\left(\tr\Phi\Phi^{\dag}\right)^{\!2} \!+\lambda_2\tr\!\left(\Phi\Phi^{\dag}\right)^{\!2} \!- \!\lambda_3\left(\det\!\Phi + \text{c.c.}\right)\!\Big]\!\Big\},
\end{gathered}
\end{align}
where $\Phi(x)$ is an $N_f$-by-$N_f$ complex matrix field, with the symmetry~\eqref{eq:3D_IR_symmetry} acting on $\Phi(x)$ in the same way as $\phi(x)$.
The question is whether this model hits a natural critical point as $\mu$ varies.

\footnote{One may speculate that the model~\eqref{eq:Landau-Ginzburg} flows to a fixed point with emergent $U(1)_A$ symmetry, with the $\lambda_3$ term irrelevant~\cite{Pisarski:1983ms,Cohen:1996ng,Lee:1996zy,Cohen:2002st,Fukushima:2008wg,Aoki:2012yj,Pelissetto:2013hqa,Cossu:2013uua}.
However, as discussed in the Supplemental Material~\ref{sec:U1A}, no convincing evidence supports this possibility in the model~\eqref{eq:Landau-Ginzburg}.
}


When $N_f=2$, the model~\eqref{eq:Landau-Ginzburg} at mean-field level undergoes an $O(4)$ second-order transition
at $\mu=|\lambda_3|$.
This critical point survive fluctuations and flows to the $O(4)$ Wilson-Fisher universality class~\cite{Pisarski:1983ms,Butti:2003nu}, as expected since the symmetry~\eqref{eq:3D_IR_symmetry} is $O(4)\rtimes(\Z_2)_{\calR}$ at $\theta_B=0,\pi$, and $SO(4)\rtimes(\Z_2)_{\calR}$ for $\theta_B\neq0,\pi$.
Hence the Landau and Landau-DQCP scenarios are both plausible for $N_f=2$.

At mean-field level, the $\lambda_3$ term drives the transition first-order for $N_f=3$, but allows an $O(2N_f^2)$ second-order transition at $\mu=0$ for $N_f\ge4$.
Unfortunately, beyond mean-field level, extensive searches have found no stable fixed points in the model~\eqref{eq:Landau-Ginzburg} for all $N_f\ge 3$~\cite{Pisarski:1983ms,Butti:2003nu,Calabrese:2004uk,Adzhemyan:2021sug,Fejos:2014qga,Resch:2017vjs,Fejos:2022mso,Fejos:2024bgl}.
We thus consider the Landau and Landau-DQCP scenarios to be disfavored for $N_f\geq 3$.




{\bf Conformal manifold.}


The only option not yet disfavored for $N_f\ge3$ is the
conformal-manifold scenario, which might work for $N_f=2$ as well. We
denote the CFT at $\theta_B$ on the critical line by
$SU(N_f)_{\theta_B}$. We obtain a non-conformal field theory by adding
the deformation
\begin{equation}
   \mu \int\d^3x\ \calO_E\,,
\end{equation}
to the $SU(N_f)_{\theta_{B}}$ CFT. This $SU(N_f)_{\theta_B}$ field theory must
describe the 3D IR physics of finite-$T$ massless QCD for all $T$. It should
flow to a gapped phase for $\mu>0$ and a Nambu-Goldstone
phase for $\mu<0$.

Whether the $SU(N_f)_{\theta_B}$ CFT exists is a nontrivial question, since
non-supersymmetric conformal manifolds are rare in
3D~\cite{Gaberdiel:2008fn,Bashmakov:2017rko,Behan:2017mwi,Hollands:2017chb,Sen:2017gfr,Giambrone:2021wsm}.
We now argue for the existence of the $SU(N_f)_0$ CFT that generates the
conformal manifold via $\calO_B$ deformations.

If we can arbitrarily ``contract'' the $\mu<0$ RG flow into its IR end by smoothly
deforming the $SU(N_f)_0$ field theory, we can study a deformed version of the $SU(N_f)_0$ CFT via
loop expansions around the Gaussian IR fixed point. The IR tail of
the $\mu<0$ RG flow is described by the 3D $SU(N_f)$ principal chiral model
(PCM) of Eq.~\eqref{eq:lowT_EFT}.
Conveniently, loop expansions in the $(2+\epsilon)$D $SU(N_f)$ PCM with
arbitrarily small $\epsilon>0$ are known to exhibit an arbitrarily short
RG flow from an interacting ultraviolet (UV) fixed point to the Gaussian
IR fixed point~\cite{ZinnJustin_book}. It is tempting to view this UV
fixed point as a weakly-coupled avatar of the $SU(N_f)_0$ CFT.

The $(2+\epsilon)$D UV fixed point should contain the avatar of the
exactly marginal operator $\calO_B$.
The key observation~\cite{Nahum:2015jya,Jones:2024ept,DeCesare:2025ukl,DeCesare:2026dwm}
is that a $d$-dimensional QFT with $d\in\R_+$ is like an
infinite-dimensional QFT, since it is formally an analytic continuation
of a family of $n$-dimensional QFTs for all $n\in\Z_+$. The
$(2+\epsilon)$D $SU(N_f)$ PCM thus sees the topology of the $SU(N_f)$
target space in all finite dimensions, including $\pi_3(SU(N_f))=\Z$.
As a result, at the $(2+\epsilon)$D UV fixed point, the closed Skyrme
3-form has protected scaling dimension exactly $3$ and generates a short
conformal multiplet~\cite{DeCesare:2025ukl}. The $(2+\epsilon)$D Skyrme 3-form is thus a
promising avatar for $\calO_B$.

A large-$N$ analysis~\cite{DeCesare:2025ukl} indicates that if the 3D
$SU(2)_0$ CFT exists, it might be close to the 3D $O(4)$ Wilson-Fisher
CFT, so existing lattice measurements~\cite{Philipsen:2021qji} might not
distinguish them. The idea that fixed points in $2+\epsilon$ and $4-\epsilon$
dimensions could have different extrapolations to 3D has recently been
highlighted by De Cesare and Rychkov~\cite{DeCesare:2025ukl,DeCesare:2026dwm}.

To see whether the $(2+\epsilon)$D UV fixed point extrapolates to
$\epsilon=1$, we study the scaling dimensions of $\calO_E$ and the lowest-dimension
primary scalar $\calO_{\phi}$ in the
$\boldsymbol{N_f}\otimes\overline{\boldsymbol{N_f}}$ representation using
the three-loop results of
Refs.~\cite{McKane:1979cm,Hikami:1980hi,Hikami:1982ak} and the four-loop results of Ref.~\cite{Wegner:1989ss}:
\begin{subequations}
\begin{align}
    \Delta_E &= 2 - \frac{1}{2}\epsilon^2 - \frac{1}{4}\epsilon^3
    - \left( \frac{5}{16} + \frac{9\zeta_3}{4N_f^2} \right)\epsilon^4
    + \calO(\epsilon^5),\\
    \Delta_{\phi} &= \left(1-\frac{1}{N_f^2}\right)\left[\epsilon
    - \frac{1}{2}  \epsilon^2 +  \frac{1}{2} \epsilon^3\right]
    +\mathcal{O}\left(\epsilon ^4\right).
\end{align}
\end{subequations}
We estimate $\Delta_E$ and $\Delta_{\phi}$ in the 3D $SU(N_f)_0$ CFT using
Pad\'e-Borel resummation; see the Supplemental Material~\ref{sec:epsilon_extrapolation}.
We summarize our
3D estimates with small $N_f$ in Table~\ref{tab:Delta_phi_estimates}, which are consistent with unitarity $\Delta_E,\Delta_{\phi} \ge \frac{1}{2}$ and relevance $\Delta_E,\Delta_{\phi} < 3$.
More terms in the $(2+\epsilon)$D expansion are needed to reduce the extrapolation errors, which are significant for $\Delta_{E}$.
\begin{table}[ht]
    \centering
    \renewcommand{\arraystretch}{1.5}
    \begin{tabular}{|c|c|c|c|c|}
        \hline
        $N_f$ & $2$ & $3$ & $4$ & $5$ \\
        \hline
        $\Delta_E$ & $1.94\pm 0.56$ & $1.76\pm 0.67$ & $1.64\pm 0.71$ & $1.56\pm 0.72$ \\
        \hline
        $\Delta_{\phi}$ & $0.54 \pm 0.02$ & $0.64 \pm 0.02$ & $0.67 \pm
        0.02$ & $0.69 \pm
        0.02$ \\
        \hline
    \end{tabular}
    \caption{Estimates for $\Delta_E$ and $\Delta_{\phi}$ in the 3D $SU(N_f)_0$ CFT obtained from Pad\'e-Borel resummations of the known $(2+\epsilon)$D expansions.}
    \label{tab:Delta_phi_estimates}
\end{table}

Conformal bootstrap~\cite{Poland:2018epd} studies exist for 3D CFTs with symmetry
$[SU(N_f)_L\times SU(N_f)_R]/\Z_{N_f}$. For $N_f = 2$,
Ref.~\cite{Kos:2015mba} found an island with $1.64\lesssim \Delta_E\lesssim 1.68$ and
$0.517 \lesssim \Delta_{\phi} \lesssim 0.521$. For $N_f=3$,
Ref.~\cite{Kousvos:2022ewl} found a rough kink at $\Delta_{\phi} = 0.62 \pm
0.01$. We see preliminary agreement at $N_f=2$
and $N_f = 3$ between our $(2 + \epsilon)$D estimates
and the bootstrap estimates.

Gathering all the evidence, we consider it promising to conjecture

\footnote{
The most ambitious conjecture would be that the 3D $SU(N_f)$ PCM~\eqref{eq:lowT_EFT} is asymptotically safe~\cite{Wilson:1973jj,Weinberg:1976xy} and gives a path-integral description of the 3D $SU(N_f)_{\theta_B}$ field theory, such that (schematically) $g_c^{-2}-g^{-2}\sim\mu$ for a critical coupling $g_c$.
}

the existence of 3D $SU(N_f)_{\theta_B}$ CFTs, though further tests are
needed. Should these CFTs exist, it would be natural for QCD to realize
the conformal-manifold scenario in the chiral limit from $N_f=2$ to the
lower edge of the conformal window.



{\bf Outlook.}


Recent lattice results~\cite{Cuteri:2021ikv,Dini:2021hug,Nakamura:2022abk,Zhang:2022kzb,Zhang:2024ldl,Zhang:2025vns,Klinger:2025xxb,DAmbrosio:2025ldv,Klinger:2026pbe}
on the QCD chiral transition motivate our conclusion that only three scenarios
(Fig.~\ref{fig:allowed_scenarios}) are consistent with a natural
second-order critical line.  Lattice QCD simulations should distinguish
them by measuring the leading IR behaviors of $\rho_B$ at
$\theta_B =0,\pi$ on the critical line, in
Table~\ref{tab:fingerprints}.
\begin{table}[ht]
    \centering
    \renewcommand{\arraystretch}{1.5}
    \begin{tabular}{|c|c|c|c|}
        \hline
         $\theta_B$ & \textbf{Landau} & \textbf{Landau-DQCP} & \textbf{Conformal-manifold}  \\
        \hline
        $0$ & irrelevant & irrelevant & exactly marginal \\ \hline
        $\pi$ & $\langle\rho_B\rangle\neq0$ & relevant & exactly marginal \\ \hline
    \end{tabular}
    \caption{The leading IR behaviors of the operator $\rho_B$ at $\theta_B =0,\pi$ on the critical line.}
    \label{tab:fingerprints}
\end{table}

The conformal-manifold scenario is the only option not disfavored for
$N_f\ge3$ by current evidence. However, it requires the
$SU(N_f)_{\theta_B}$ universality class, whose existence is an open
question of independent interest. Demonstrating it will require improved
conformal bootstrap bounds on 3D CFTs with the
symmetry~\eqref{eq:3D_IR_symmetry} and further development of the
$(2+\epsilon)$D expansion. Lattice Monte Carlo simulations of the 3D
$SU(N_f)$ PCM using L\"uscher's admissibility
method~\cite{Luscher:1981zq,Schramm:2000cc,Shiozaki:2024yrm,Lopez-Contreras:2025iik}
(see also Ref.~\cite{Chen:2024ddr}) might be valuable.

While we assume a second-order critical line, a \emph{weakly}
first-order transition is not excluded by
Refs.~\cite{Cuteri:2021ikv,Dini:2021hug,Nakamura:2022abk,Zhang:2022kzb,Zhang:2024ldl,Zhang:2025vns,Klinger:2025xxb,DAmbrosio:2025ldv,Klinger:2026pbe}.
If such transitions are controlled by nearby
non-unitary complex CFTs~\cite{Gorbenko:2018ncu,Gorbenko:2018dtm}, such
as the $(2+\epsilon)$D fixed points, much of our discussion would carry
over.

Finally, we should extend our analysis to real chemical
potentials. Thermal QCD with real $\mu_B$ violates 3D reflection
positivity (i.e.~unitarity) but remains ``real'' thanks to anti-unitary
symmetries related to time reversal. The $SU(N_f)_{\theta_B}$
conformal manifold might analytically continue to a strip around
$\theta_B\in\R$. If so, the chiral phase transition could be described
by non-unitary real CFTs~\cite{Gorbenko:2018ncu,Gorbenko:2018dtm} for
small real $\mu_B$, and become first-order for larger real $\mu_B$.
Testing this requires further work.



{\bf Acknowledgments.}


We are very grateful to Thomas Dumitrescu, Amey Gaikwad, Evgenii Ievlev, Theodore Jacobson, David B. Kaplan, Joseph Kapusta, Zohar Komargodski, Owe Philipsen, Jiaxin Qiao, Srimoyee Sen, Laurence Yaffe, and Yuan Xin for very helpful discussions.
OpenAI's Codex
5.4, Anthropic's Claude Opus 4.6, and Refine.ink were helpful in
proofreading the manuscript. A.C. and S.C. are supported by Simons
Foundation award number 994302, as part of the Simons Collaboration on Confinement and QCD Strings. R.D.P. is supported by the U.S.
Department of Energy under contract DE-SC0012704, and by the Alexander
von Humboldt Foundation.

\bibliography{non_inv}

\onecolumngrid
\appendix
\renewcommand{\thesection}{S\arabic{section}}
\numberwithin{equation}{section}

\newpage

\section*{Supplemental Material}




\section[Effective theory in the high-T gapped phase]{Effective theory in the high-$T$ gapped phase}
\label{sec:high_T_EFT}




According to perturbative calculations~\cite{KapustaGale201102, Roberge:1986mm}, the high-$T$ (i.e.~small-$\beta$) phase of massless QCD is gapped for all $\theta_B$.
The free energy density $f_{\beta}(\theta_B)$ of Eq.~\eqref{eq:free_energy} at small $\beta$ is an even $2\pi$-periodic analytic function for every $\theta_B$ except at $\theta_B=\pi$ (mod $2\pi$), where it has a cusp-like maximum.
All the propagating degrees of freedom are gapped out at distance scales $\ell \gg \beta$, so a 3D effective field theory at $\ell \gg \beta$ describes the colorless vacua only.
Such a 3D QFT is not completely trivial because the theory is not uniformly trivially gapped for all $\theta_B$.
It has a path-integral description in terms of two colorless fields, a real scalar field $\varphi$ and a 2-form $U(1)$ gauge field $b$ \cite{Kallosh:1995hi,Gabadadze:1999na,Gabadadze:2000vw,Gabadadze:2002ff,Dvali:2005an,Burgess:2023ifd}:
\begin{equation}\label{eq:high_T_EFT}
    \int\calD b\calD\varphi \exp\left\{ -\int\d^3x\,
    \bar{f}_{\beta}(\varphi) + \frac{i}{2\pi} \int\varphi\d b\right\}\,.
\end{equation}
The above $\bar{f}_{\beta}(\varphi)$ is an even analytic function of $\varphi$ such that
\begin{equation}
    \bar{f}_{\beta}(\varphi)
    \begin{cases}
        = f_{\beta}(\varphi), & -\pi\le\varphi\le\pi \\
        > f_{\beta}(\varphi), & \text{for other }\varphi
    \end{cases}.
\end{equation}
This 3D model has a $2$-form $U(1)$ symmetry on $b$'s electric side and a $(-1)$-form $U(1)$ symmetry on $b$'s magnetic side.
With a background $3$-form $U(1)$ gauge field $C$ and a background $0$-form $U(1)$ gauge field (i.e.~a background $2\pi$-periodic scalar) $\theta_B$, the path integral reads
\begin{equation}\label{eq:high_T_EFT_background}
    \int\calD b\calD\varphi \exp\left\{ -\int\d^3x\,
    \bar{f}_{\beta}(\varphi) + \frac{i}{2\pi} \int (\varphi - \theta_B)(\d b-C)\right\}\,.
\end{equation}
This 3D model has no propagating degrees of freedom and describes infinitely many vacua (sometimes called ``universes''~\cite{Hellerman:2006zs,Sharpe:2014tca,Tanizaki:2019rbk,Komargodski:2020mxz,Cherman:2020cvw,Lin:2025oml,Sharpe:2022ene}) labeled by the $2$-form $U(1)$ symmetry charge $n\in\Z$.
For constant $\theta_B$, each vacuum has the energy density given by $\bar{f}_{\beta}(2\pi n + \theta_B)$, and two vacua cross at $\theta_B=\pi$.
The $2$-form $U(1)$ symmetry has to be emergent and only the lowest vacua are relevant for high-$T$ massless QCD.

To match the 3D 't Hooft anomaly~\eqref{eq:SPT_3D}, the high-$T$ effective theory~\eqref{eq:high_T_EFT} is supposed to see the full 3D symmetry~\eqref{eq:3D_IR_symmetry}.
The actions of discrete symmetries on the high-$T$ effective theory~\eqref{eq:high_T_EFT} are evident:
\begin{equation}
    \varphi(x)\xrightarrow{\calC}-\varphi(x)\,,\
    b(x)\xrightarrow{\calC}-b(x)\,,\qquad
    \varphi(x)\xrightarrow{\calS}-\varphi(x)\,,\
    b(x)\xrightarrow{\calS}-b(x)\,,\qquad
    \varphi(x)\xrightarrow{\calR}\varphi(\overline{x})\,,\
    b(x)\xrightarrow{\calR}-\overline{b}(\overline{x})\,.
\end{equation}
At first glance, it seems that there is no room for the $SU(N_f)_{L,R}$ symmetry to act, since Eq.~\eqref{eq:high_T_EFT_background} has no continuous 0-form symmetry of any sort.
Nevertheless, the emergent $2$-form $U(1)$ symmetry can be a homomorphism target of the $0$-form $SU(N_f)_{L,R}$ symmetry.
Namely, the $A_{L,R}$ gauge fields can be coupled via the substitution
\begin{equation}\label{eq:C=A^2}
    C = \frac{1}{4\pi}\tr\!\left( A_L \d A_L - A_R \d A_R
    - \i\frac{2}{3} A_L^3 + \i\frac{2}{3} A_R^3 \right).
\end{equation}
The emergent $2$-form $U(1)$ symmetry has a mixed 't Hooft anomaly with the $(-1)$-form $U(1)$ symmetry, described by the 4D invertible theory
\begin{equation}\label{eq:high_T_anomaly}
    \exp\left\{\frac{i}{2\pi}\int\theta_B\d C\right\}.
\end{equation}
This emergent anomaly precisely reproduces the 3D anomaly~\eqref{eq:SPT_3D} via the substitution~\eqref{eq:C=A^2}.
Mathematically, such a homomorphism from $0$-form $SU(N_f)_{L,R}$ to $2$-form $U(1)$ is encoded in the classifying-space fibration $BSU(N_f)\to B^3U(1)$ that induces an isomorphism from $\pi_4(BSU(N_f))=\Z$ to $\pi_4(B^3U(1))=\Z$.

The above way of matching anomalies via flowing to an emergent higher-form symmetry was recently packaged into the notion of symmetry transmutation~\cite{Seiberg:2025bqy}.
So we can say that the $0$-form $SU(N_f)_{L,R}$ chiral symmetry of QCD is transmuted into a $2$-form $U(1)$ symmetry in its high-$T$ effective theory.





\section[Second-order phase transition at theta=pi in the high-T phase]{Second-order phase transition at $\theta_B=\pi$ in the high-$T$ phase}
\label{sec:anomalons}




In the high-$T$ phase, a window of second-order phase transitions at $\theta_B=\pi$ would have a chance to be realized at relatively lower $T$, if a 3D CFT consistent with symmetry and anomaly existed.
The upper edge of this gapless window would be a tricritical point, while the lower edge could be connected to the DQCP or the $SU(N_f)_{\pi}$ CFT on the critical line of chiral transitions.

One may speculate that such an appropriate 3D CFT can simply be provided by colorless free fermions, because the 3D anomaly~\eqref{eq:SPT_3D} can be matched by an $N_f$-flavor free massless 4D Dirac fermion at finite temperatures.
At $\theta_B=\pi$, one Matsubara mode becomes massless, leading to the 3D effective theory at $\ell\gg\beta$:
\begin{align}\label{eq:anomalon}
    \int\calD\psi\calD\bar{\psi}\calD\chi\calD\bar{\chi}
    \ \exp\left\{-\int d^{3}x\, \Big(\bar{\psi}\slashed{\partial}\psi +\bar{\chi}\slashed{\partial}\chi\Big)\right\}\,,
    \qquad(\theta_B=\pi)
\end{align}
with two $N_f$-flavor free massless 3D Dirac fermions $\psi$ and $\chi$.
At $\theta_B\neq\pi$, all the Matsubara modes are massive and these 3D fermions are gapped out at $\ell\gg\beta$.
This speculation has an interplay with the Landau-DQCP scenario.
Around the chiral transition, these 3D fermions have to talk to the chiral condensate $\Phi$ through Yukawa couplings,
\begin{equation}
    y_1\int d^3x \left(\bar{\psi} \Phi \Phi^{\dag} \psi - \bar{\chi} \Phi^{\dag} \Phi \chi\right)\,,
    \qquad
    y_2\int d^3x \left(\bar{\psi}\Phi\chi + \bar{\chi}\Phi^{\dag}\psi \right)\,,
    \qquad\text{etc.}
\end{equation}
Therefore, the DQCP at $\theta_B=\pi$ might be described by a certain 3D Gross-Neveu-Yukawa universality class, as long as the gapless window exists and has a free-fermion description.

Nevertheless, as already pointed out by Ref.~\cite{Kobayashi:2023ajk}, although it matches the anomaly~\eqref{eq:SPT_3D} of QCD, the free-fermion CFT~\eqref{eq:anomalon} does not match the symmetry~\eqref{eq:QCD_symmetry} of QCD.
A colorless fermionic object in the fundamental representation of $SU(N_f)_L$ or in the anti-fundamental representation of $SU(N_f)_R$ can exist in QCD only when
\begin{equation}
    N_c = 1 \mod N_f
    \qquad\text{and}\qquad
    N_c = 1 \mod 2
\end{equation}
When this happens, this colorless fermionic object necessarily carries a baryon number $q_B$ satisfying $q_B = 1 \mod N_f$.
For example, with physical $N_c=3$, the above conditions are satisfied by $N_f=2$, but are violated by $N_f=3$.
Therefore, a gapless window at $\theta_B=\pi$ described by the free-fermion CFT~\eqref{eq:anomalon} is at least conceivable for $N_f=2$, but is immediately excluded for $N_f=3$.
Of course, a more complicated free-fermion CFT might simultaneously match the symmetry and the anomaly, and we leave a detailed exploration of such possibilities to the future.




\section[The fate of emergent U(1)A symmetry at the chiral transition]{The fate of emergent $U(1)_A$ symmetry at the chiral transition}
\label{sec:U1A}




It is conceivable that a $U(1)_A$ symmetry emerges on the critical line in Landau scenarios~\cite{Cohen:1996ng,Lee:1996zy,Cohen:2002st,Fukushima:2008wg,Aoki:2012yj,Aoki:2012yj,Pelissetto:2013hqa,Cossu:2013uua,Pisarski:2024esv}.
For this to happen, two requirements are needed: (I) the Ginzburg-Landau model~\eqref{eq:Landau-Ginzburg} with $\lambda_3=0$ must have an IR fixed point and (II) $\det\Phi$ must be an irrelevant operator at that fixed point.
Establishing (I) is already highly nontrivial since the $\lambda_2$ fluctuations destabilize the mean-field $O(2N_f^2)$ critical point.

When $N_f=2$, the model~\eqref{eq:Landau-Ginzburg} with $\lambda_3=0$ may flow to some Wilson-Fisher-like $O(4)\times O(2)$ fixed point~\cite{Pelissetto:2013hqa,Nakayama:2014sba,Henriksson:2020fqi} (see also~Refs.~\cite{Yabunaka:2017uox,Yabunaka:2018mju,Sorokin:2021jwf,Sorokin:2022zwh,Sanchez-Villalobos:2024vmd}).
But $\det\Phi\sim\Phi^2$ is relevant at that fixed point, since it
is literally a mass term.
When $N_f\ge3$, there is no convincing evidence for the existence of fixed points in the model~\eqref{eq:Landau-Ginzburg} with $\lambda_3=0$~\cite{Pisarski:1983ms,Butti:2003nu,Calabrese:2004uk,Adzhemyan:2021sug,Fejos:2014qga,Resch:2017vjs,Fejos:2022mso,Fejos:2024bgl}.
We thus regard a second-order chiral transition with emergent $U(1)_A$ as disfavored in massless QCD for all $N_f\ge 2$ given current evidence.

Further searches for fixed points in the model~\eqref{eq:Landau-Ginzburg} with $\lambda_3=0$ using lattice Monte Carlo and other methods would still be valuable.
This is because if such a fixed point existed, $\det\Phi\sim\Phi^{N_f}$ should be irrelevant for sufficiently large $N_f$, and thus $U(1)_A$ would have a chance to emerge at QCD chiral transition with sufficiently large $N_f$.
Nevertheless, it is hard to imagine that anything like this can happen for $N_f=3$, since it seems very difficult to get $\det\Phi\sim\Phi^3$ to be irrelevant.




\section[Extrapolating from (2+epsilon)D to 3D]{Extrapolating from $(2+\epsilon)$D to $3$D}
\label{sec:epsilon_extrapolation}



The beta function of the $(2+\epsilon)$D $SU(N_f)$ PCM is known up to the four-loop order~\cite{Hikami:1980hi,Hikami:1982ak,Wegner:1989ss} (their coupling $t$ in the Lagrangian is related to our $g^2$ through $g^2=4t$):
\begin{align}
\begin{split}
    \beta(g^2) \ =&\: \ \epsilon g^2 -\frac{1}{4} N_f g^4
    -\frac{1}{32} N_f^2 g^6-\frac{3}{512} N_f^3 g^8 - \frac{1}{1024}\left[\frac{19}{12}N_f^4 + 3\zeta_3N_f^2 \right]g^{10}
    + \calO\left(g^{12}\right)\,,
\end{split}
\end{align}
where $\zeta_3\approx 1.202$ is the Riemann zeta function evaluated at $3$.
The UV fixed point is thus located at
\begin{equation}
    g_*^2\ =\ \frac{4}{N_f}\epsilon - \frac{2}{N_f}\epsilon^2 + \frac{1}{2N_f}\epsilon^3
    - \left(\frac{1}{3N_f}
    +\frac{3\zeta_3}{N_f^3}\right)\epsilon^4
    + \calO(\epsilon^5)
\end{equation}
The scaling dimension of the energy operator $\calO_E$ is related to the critical exponent $\nu$ that governs the divergence of the correlation length as the UV fixed point:
\begin{equation}\label{eq:Delta_energy}
\begin{split}
    \Delta_E(\epsilon)
    \ =\
    2+\epsilon - \frac{1}{\nu}
    \ =\
    2 + \epsilon +\beta'(g^2_{*}) \ =\
    2 - \frac{1}{2}\epsilon^2 - \frac{1}{4}\epsilon^3
    - \left( \frac{5}{16} + \frac{9\zeta_3}{4N_f^2} \right)\epsilon^4
    + \calO(\epsilon^5)
    \,.
\end{split}
\end{equation}
This series appears to be non-alternating. Curiously, $N_f$-dependence only
shows up from the fourth order.

The anomalous dimension of the vertex operator $\calO_{\phi}$ in the $\boldsymbol{N_f}\otimes\overline{\boldsymbol{N_f}}$ representation of $SU(N_f)_L\times SU(N_f)_R$ is known up to the three-loop order~\cite{McKane:1979cm} (their coupling $g^2$ in the Lagrangian is the same as ours):
\begin{align}
    \gamma_{\phi}(g^2) \  = \ \frac{\left(N_f^2-1\right)}{4 N_f} \left(g^2
    + \frac{3 N_f^2}{128}  g^6\right) +\calO\left(g^8\right)\,.
\end{align}
This $\gamma_{\phi}$ is the anomalous dimension with respect to the engineering dimension $0$.
Therefore, this $\gamma_{\phi}$ is equivalent to $\frac{1}{2}\zeta$ used in Ref.~\cite[Section~4.1.3]{Henriksson:2022rnm} (see especially their Footnote~35) for the $O(N)$ nonlinear sigma model.
Hence the scaling dimension of $\phi$ at the UV fixed point is given by
\begin{align}
    \label{eq:Delta_phi}
    \Delta_{\phi}(\epsilon)
    \ =\ \gamma_{\phi}(g_*^2)
    \ = \ \left(1-\frac{1}{N_f^2}\right)\left[\epsilon
    - \frac{1}{2}  \epsilon^2 +  \frac{1}{2} \epsilon^3\right]
    +\calO\left(\epsilon ^4\right)  \,.
\end{align}
This series appears to be alternating.
The factorization of the $N_f$-dependence is unlikely to hold from the fourth order.

To extrapolate $\Delta_E(\epsilon)$ and $\Delta_{\phi}(\epsilon)$ to $\epsilon=1$, we need to deal with a standard challenge.
Since there are no non-trivial fixed points for $\epsilon<0$, Dyson's
argument~\cite{Dyson:1952tj} suggests that the perturbative expansion in $\epsilon$
should be asymptotic.  Simply plugging $\epsilon = 1$ into a truncated series representation of $\Delta_E(\epsilon)$ and $\Delta_{\phi}(\epsilon)$ is thus likely to be a very poor approximation.  However, it is natural to expect observables like
$\Delta_{E}(\epsilon), \Delta_{\phi}(\epsilon)$ to have resurgent transseries
representations around $\epsilon = 0$, see e.g.~Ref.~\cite{Aniceto:2018bis} for
a review.  The standard way to extrapolate asymptotic series is via Borel
summation of an analytic continuation of the Borel transform of the series. Here
we will adopt the
simplest and most widely used analytic continuation, which is via Pad\'e
approximants, which are rational functions that match the series expansion to a
given order.  We will use this Pad\'e-Borel (PB) resummation approach to estimate the values of $\Delta_E(\epsilon)$ and $\Delta_{\phi}(\epsilon)$ at $\epsilon=1$.  To estimate the errors on our estimates, we will use combinations of the variation of our results as we vary the choice of Pad\'e approximant together  with the non-perturbative ambiguities in the resummations (when such ambiguities appear).





\subsection[Extrapolation of Delta-energy]{Extrapolation of $\Delta_E$}
\label{sec:delta_energy}



Replacing $\epsilon^n$ with $\delta^n/n!$, we obtain the Borel transform of the truncated $\Delta_E(\epsilon)$ series~\eqref{eq:Delta_energy}:
\begin{align}\label{eq:Borel_energy}
    \calB[\Delta_E](\delta)\ = 2-\frac{\delta^2}{4}
    - \frac{\delta^3}{24}
    - \left( \frac{5}{16} + \frac{9\zeta_3}{4N_f^2} \right)\frac{\delta^4}{24}
    +\calO\left(\delta^5\right)\,.
\end{align}
If $\Delta_E(\epsilon)$ is indeed non-alternating, we would expect $\calB[\Delta_E](\delta)$ to have a leading singularity on the positive real axis.
Such a singularity leads to imaginary ambiguities in the Borel integral that must be canceled by appropriate nonperturbative contributions in a resurgent transseries representation of $\Delta_E(\epsilon)$.

An $[m/n]$ Pad\'e approximant to $\calB[\Delta_E](\delta)$ is a degree-$(m,n)$ rational function that matches the series expansion of $\calB[\Delta_E](\delta)$ to order $\delta^{m+n}$.
The near-diagonal Pad\'e approximants at order $\delta^4$ are
\begin{subequations}
\begin{align}
    \calB[\Delta_E](\delta)\big|_{[3/1]} &=
    \frac{2 - \Bigl(\frac{5}{8}\!+\!\frac{9\zeta_3}{2N_f^2}\Bigr)\delta - \frac{1}{4}\delta^2 + \Bigl(\frac{7}{192}\!+\!\frac{9\zeta_3}{16N_f^2}\Bigr)\delta^3 }
    {1 - \Bigl(\frac{5}{16}\!+\!\frac{9\zeta_3}{4N_f^2}\Bigr)\delta}\,,\\
    \calB[\Delta_E](\delta)\big|_{[2/2]} &=
    \frac{2 - \frac{1}{3}\delta - \Bigl(\frac{43}{144}\!+\!\frac{3\zeta_3}{4N_f^2}\Bigr)\delta^2}
    {1 - \frac{1}{6}\delta - \Bigl(\frac{7}{288}\!+\!\frac{3\zeta_3}{8N_f^2}\Bigr)\delta^2}\,,\\
    \calB[\Delta_E](\delta)\big|_{[1/3]} &=
    \frac{2 - \Bigl(\frac{17}{8}\!+\!\frac{9\zeta_3}{2N_f^2}\Bigr)\delta }
    {1 - \Bigl(\frac{17}{16}\!+\!\frac{9\zeta_3}{4N_f^2}\Bigr)\delta + \frac{1}{8}\delta^2 - \Bigl(\frac{43}{384}\!+\!\frac{9\zeta_3}{32N_f^2}\Bigr)\delta^3}\,.
\end{align}
\end{subequations}
All of them have poles on the positive real axis.
We then obtain the $[m/n]$ PB-resummed function $\Delta_E(\epsilon)$ by performing the Borel integral,
\begin{align}
    \Delta_E(\epsilon)\big|_{[m/n]} =\  \lim_{\theta\to0}\  \frac{1}{\epsilon}\int_0^{\infty} d \delta\, \exp\left(-\delta e^{i \theta}/\epsilon \right)\, \calB[\Delta_E]( \delta e^{i \theta} )\big|_{[m/n]}\,.
\end{align}
In a resurgent transseries representation of $\Delta_E(\epsilon)$, $\im \Delta_E(\epsilon)\big|_{[m/n]}$ would be canceled by a corresponding imaginary part of a non-perturbative term in the transseries, which would also give a real contribution of approximately the same magnitude.  We thus use the magnitude of the imaginary ambiguities in the Pad\'e-Borel resummations as an estimate of the error on our extrapolations.


We show the functions $\Delta_E(\epsilon)|_{[m/n]}$ with $N_f=2\sim7$ in Fig.~\ref{fig:Pade_Borel_E}.
We summarize our $\epsilon=1$ results as a function of $[m/n]$ and $N_f$ in Table~\ref{tab:Pade_sums_E}.
We use the union of the envelopes of $\Delta_E(1)|_{[3/1]}$, $\Delta_E(1)|_{[2/2]}$, and $\Delta_E(1)|_{[1/3]}$ to determine our estimation of 3D $\Delta_E$ and the extrapolation error, as we summarize in Table~\ref{tab:Pade_sums_E'}.
The extrapolation errors are large with the currently available four-loop data, and thus higher-loop data would be very valuable.

\begin{figure*}[ht]
  \centering
  \includegraphics[width=0.5\textwidth]{figures/E_borel_pade_nf2.pdf}
  \hfill
  \includegraphics[width=0.5\textwidth]{figures/E_borel_pade_nf3.pdf}
  \hfill
  \\
  \includegraphics[width=0.5\textwidth]{figures/E_borel_pade_nf4.pdf}
  \hfill
  \includegraphics[width=0.5\textwidth]{figures/E_borel_pade_nf5.pdf}
  \hfill
  \\
  \includegraphics[width=0.5\textwidth]{figures/E_borel_pade_nf6.pdf}
  \hfill
  \includegraphics[width=0.5\textwidth]{figures/E_borel_pade_nf7.pdf}

  \caption{Evaluations of the PB-resummed function $\Delta_E(\epsilon)$ with $N_f=2\sim7$. The solid lines show $\mathrm{Re}\,\Delta_E(\epsilon)$ while the shadow bands indicate error bands determined from the imaginary ambiguities in PB resummation.}
  \label{fig:Pade_Borel_E}
\end{figure*}

\begin{table}[ht]
    \centering
    \renewcommand{\arraystretch}{1.5}
    \begin{tabular}{|c|c|c|c|c|}
        \hline
        $[m/n]$ & $[3/1]$ & $[2/2]$ & $[1/3]$ \\
        \hline
        $\Delta_E(1),\ {N_f=2}$ & $1.64\pm \,0.05$ & $1.94\pm \,0.56$ & $1.76\pm \,0.01$ \\
        \hline
        $\Delta_E(1),\  {N_f=3}$ & $1.71\pm \,0.18$ & $1.76\pm \,0.67$ & $1.77\pm \,0.01$ \\
        \hline
        $\Delta_E(1),\  {N_f=4}$ & $1.72\pm \,0.31$ & $1.64\pm \,0.71$ & $1.77\pm \,0.02$ \\
        \hline
        $\Delta_E(1),\  {N_f=5}$ & $1.69\pm \,0.39$ & $1.56\pm \,0.72$ & $1.78\pm \,0.03$ \\
        \hline
        $\Delta_E(1),\  {N_f=6}$ & $1.66\pm \,0.44$ & $1.51\pm \,0.72$ & $1.78\pm \,0.03$ \\
        \hline
        $\Delta_E(1),\  {N_f=7}$ & $1.64\pm \,0.47$ & $1.48\pm \,0.72$ & $1.78\pm \,0.03$ \\
        \hline
    \end{tabular}
    \caption{Evaluations of PB-resummed $\Delta_E(1)$ with $N_f=2\sim7$. }
    \label{tab:Pade_sums_E}
\end{table}

\begin{table}[ht]
    \centering
    \renewcommand{\arraystretch}{1.5}
    \begin{tabular}{|c|c|c|c|c|c|c|}
        \hline
        $N_f$ & $2$ & $3$ & $4$ & $5$ & $6$ & $7$\\
        \hline
        $\Delta_E$ & $1.94\pm 0.56$ & $1.76\pm 0.67$ & $1.64\pm 0.71$ & $1.56\pm 0.72$ & $1.51\pm 0.72$ & $1.48\pm 0.72$ \\
        \hline
    \end{tabular}
    \caption{Estimates of 3D $\Delta_E$ obtained from the PB resummation of the four-loop truncated $(2+\epsilon)$D perturbative expansion, with $N_f=2\sim7$.}
    \label{tab:Pade_sums_E'}
\end{table}




\subsection[Extrapolation of Delta-phi]{Extrapolation of $\Delta_{\phi}$}
\label{sec:delta_phi}



Replacing $\epsilon^n$ with $\delta^n/n!$, we obtain the Borel transform of the truncated $\Delta_{\phi}(\epsilon)$ series~\eqref{eq:Delta_phi}:
\begin{align}
    \calB[\Delta_{\phi}](\delta)\ =\ \left(1-\frac{1}{N_f^2}\right)\left[\delta-\frac{\delta^2}{4}
    +\frac{\delta^3}{12}\right]+\calO\left(\delta^4\right) \,.
\end{align}
If $\Delta_{\phi}(\epsilon)$ is indeed alternating, we would expect $\calB[\Delta_{\phi}](\delta)$ to have a leading singularity on the negative real axis, and the Borel integral might be free of imaginary ambiguities.
The near-diagonal Pad\'e approximants up to order $\delta^3$ are
\begin{subequations}
\begin{align}
    \calB[\Delta_{\phi}](\delta)\big|_{[1/1]} &=
    \left(1-\frac{1}{N_f^2}\right)
    \frac{\delta}{1+\frac{1}{4}\delta}\,,\\
    \calB[\Delta_{\phi}](\delta)\big|_{[2/1]} &=
    \left(1-\frac{1}{N_f^2}\right)
    \frac{\delta + \frac{1}{12}\delta^2}{1+\frac{1}{3}\delta}\,,\\
    \calB[\Delta_{\phi}](\delta)\big|_{[1/2]} &=
    \left(1-\frac{1}{N_f^2}\right)
    \frac{\delta}{1+\frac{1}{4}\delta-\frac{1}{48}\delta^2}\,.
\end{align}
\end{subequations}
The $[1/1]$ and $[2/1]$ Pad\'e approximants have no singularities on the positive real axis, while the $[1/2]$ Pad\'e approximant has a pole at $\delta\approx 15.2 \gg 0$.
We then obtain the $[m/n]$ PB-resummed function $\Delta_{\phi}(\epsilon)$ by performing the Borel integral,
\begin{align}
    \Delta_{\phi}(\epsilon)\big|_{[m/n]} =  \frac{1}{\epsilon}  \int_0^{\infty} d \delta\, e^{-\delta/\epsilon}\, \calB[\Delta_{\phi}](\delta)\big|_{[m/n]}\,.
\end{align}
We show the functions $\Delta_{\phi}(\epsilon)|_{[m/n]}\big/\bigl(1-\frac{1}{N_f^2}\bigr)$ in Fig.~\ref{fig:Pade_Borel_phi}.
The solid lines show the real parts while the shadow bands indicate the imaginary ambiguities, which only appear for $[1/2]$ and are essentially invisible.  (The imaginary ambiguities only become appreciable for $\epsilon \gtrsim 2$.)
We summarize the results at $\epsilon=1$ in Table~\ref{tab:Pade_sums_phi}.
We take the union of the envelopes of $\Delta_{\phi}(1)|_{[1/1]}$, $\Delta_{\phi}(1)|_{[2/1]}$, and $\Delta_{\phi}(1)|_{[2/1]}$ as our final estimates of 3D $\Delta_{\phi}$, as we summarize in Table~\ref{tab:Pade_sums_phi'}.
We would need more terms beyond the three-loop order in the $(2+\epsilon)$D expansion to get more accurate results.

\begin{figure*}[ht]
  \centering
  \includegraphics[width=0.5\textwidth]{figures/phi_borel_pade.pdf}

  \caption{Evaluations of PB-resummed $\Delta_{\phi}(\epsilon)\big/\bigl(1-\frac{1}{N_f^2}\bigr)$. The solid lines show $\mathrm{Re}\,\Delta_{\phi}(\epsilon)\big/\bigl(1-\frac{1}{N_f^2}\bigr)$ and the shadow bands indicate $\pm\mathrm{Im}\,\Delta_{\phi}(\epsilon)\big/\bigl(1-\frac{1}{N_f^2}\bigr)$.  The imaginary ambiguities only appear for $[1/2]$ and are too small to be visible.}
  \label{fig:Pade_Borel_phi}
\end{figure*}

\begin{table}[ht]
    \centering
    \renewcommand{\arraystretch}{1.5}
    \begin{tabular}{|c|c|c|c|c|}
        \hline
        $[m/n]$ & $[1/1]$ & $[2/1]$ & $[1/2]$ \\
        \hline
        $\Delta_{\phi}(1)\big/\bigl(1-\frac{1}{N_f^2}\bigr)$ & $0.70$ & $0.73$ & $0.74\pm i\,0.00$\\
        \hline
    \end{tabular}
    \caption{Evaluations of PB-resummed $\Delta_{\phi}(1)\big/\bigl(1-\frac{1}{N_f^2}\bigr)$. }
    \label{tab:Pade_sums_phi}
\end{table}

\begin{table}[ht]
    \centering
    \renewcommand{\arraystretch}{1.5}
    \begin{tabular}{|c|c|c|c|c|c|c|}
        \hline
        $N_f$ & $2$ & $3$ & $4$ & $5$ & $6$ & $7$\\
        \hline
        $\Delta_{\phi}$ & $0.54\pm 0.02$ & $0.64\pm 0.02$ & $0.67\pm 0.02$ & $0.69\pm 0.02$ & $0.70\pm 0.02$ & $0.71\pm 0.02$ \\
        \hline
    \end{tabular}
    \caption{Estimates of 3D $\Delta_{\phi}$ obtained from the PB resummation of the three-loop truncated $(2+\epsilon)$D perturbative expansion, with $N_f=2\sim7$.}
    \label{tab:Pade_sums_phi'}
\end{table}

\end{document}
===== END FILE: main.tex =====
